% Transcription — fonds Alexandre Grothendieck, shelfmark 106, batch 1 (pages 1-16).
% Established from the facsimile archives/batches/106/batch-01.pdf (a mirror of
% https://grothendieck.umontpellier.fr/106.pdf), per the skill
% transcribe-grothendieck. The folder is sixteen pages and this is its only
% batch.
%
% Eleven of the sixteen leaves are transcribed. The folder is physically one
% thing: eight spiral-notebook leaves in ink, written recto only, interleaved
% with the sheets of a computer listing he used as scrap. His own pagination
% stands at the head of the notebook leaves and runs 1 to 8 — the numeral is
% visible on pages 3, 5, 7, 9, 11, 13 and 15 as 2) to 8), page 1 opening the
% run without a numeral — so the archivists' odd pages 1, 3, 5, 7, 9, 11, 13
% and 15 are his pages 1 to 8 and the run is continuous and complete. The even
% pages are the listing. Three of them carry mathematics of his: page 2, a
% landscape sheet turned sideways and covered end to end with trial diagrams;
% page 12, two small trial figures in the bottom margin; page 14, the block of
% amalgamated-sum squares that page 15 uses. Pages 4, 6, 8, 10 and 16 carry
% nothing of his and are skipped, the gaps in the page numbers being the
% record. Page 8 does carry pencil, but it is the machine room's — a job
% number beside the listing's own header — and not his.
%
% The inventory's title, « [Champs (stacks) 3] », is bracketed and is
% therefore the archivists'; nothing in the folder carries a title of his, and
% the two section headings used here are the edition's and are marked as such.
% What the folder is about is narrower than « champs » and is named on its
% first line: it is the localisation of a category with weak equivalences and
% cofibrations, in the axiomatics of Baues, and specifically the question
% whether the homotopy category can be computed by a calculus of fractions of
% one single kind. The neighbouring folder 107, « Closed models », is the same
% subject under its other name.
%
% The run is one argument and it fails, in the open, twice. Page 1 fixes
% (C, W, cof) satisfying C1 to C3 de Baues and asserts that
% W_cof^{-1} C_cof --> W^{-1} C is an equivalence; to prove it a quasi-inverse
% is built by cofibrant replacement — C3 applied to the initial map gives
% alpha_x : x-tilde --> x in W with x-tilde cofibrant — and a map f : x --> y
% is lifted by the « weak lifting lemma » not into y-tilde but into some
% y-tilde', so that the arrow of the quasi-inverse is the fraction
% [i]^{-1}[f-tilde]. Independence of the choices is reduced, by forming the
% amalgamated sum of two such lifts over y-tilde, to one statement: if
% x', y' are cofibrant and alpha.phi = alpha.psi with alpha in W, then
% [phi] = [psi]. Page 3 checks that the functor inverts W and that the two
% composites are the identity up to isomorphism, the second by the functorial
% identity [alpha_y]([i_f]^{-1}[f-tilde]) = [f][alpha_x] which falls out of
% alpha_y = beta_f i_f. Page 5 states the resulting Prop.: Psi is an
% equivalence if and only if the condition (*) holds, (*) being exactly the
% statement isolated on page 1, and then changes the ground — from here on
% every object is assumed cofibrant — and defines, for x and y, the category
% Ch_0(x,y) of diagrams x --> y-tilde <-- y with the backward arrow in
% cof inter W, together with a composition of those categories. Page 7 checks
% that the composition is associative up to canonical isomorphism and unital,
% passes to pi_0, and gets from it a category he writes C-tilde and a functor
% C --> C-tilde; then asks whether that functor inverts W. Page 9 applies the
% weak lifting lemma to get i in W inter cof and alpha in W with alpha.i = f
% and a section g of alpha, proposes (g,i) as the inverse, and then breaks
% off: the composite alpha.i, which is f, lives in W and not in cof, and
% « ça se présente mal ! ». So the order of business is reversed and the
% canonical functor q : C-tilde --> W^{-1} C, x |--> x, (f,i) |--> [i]^{-1}[f]
% is defined instead and shown bijective on objects and surjective on arrows.
% Page 11 is the second failure and the folder's turning point: it is not
% clear that q is faithful, and inverting only the cofibrant weak
% equivalences — W_0 = W inter cof — is doubted in as many words; the
% definition is therefore rectified, Ch_0(x,y) being taken from then on with
% f in cof and i in W rather than with i in cof inter W. Pages 13 and 15 work
% the rectified definition: a map is factored as a cofibration followed by a
% weak equivalence, two such factorisations are compared through a third built
% as an amalgamated sum, and the run closes on the criterion (**) — if f in W
% has two sections i and i', the class of (1_x, i) in pi_0(Ch_0(x,y)) must not
% depend on the section — which is stated to be necessary and sufficient for
% q to be faithful, and is not proved.
%
% The hand is drafting throughout and that is the folder's governing fact.
% Pages 1, 3, 5 and 9 read almost end to end; pages 11, 13 and 15 are a fast
% palimpsest whose statements, displayed formulas and underlined words carry
% and whose connective prose often does not, page 13 being the worst of the
% folder and carrying besides a marginal note written up the left edge at an
% angle that will not resolve at the resolution the film holds. Following
% folder 29 batches 8 and 9, folder 52, folder 41, folder 88, folder 152 and
% folder 104, what the leaves will not support is left \ill{} rather than
% invented, and the apparatus is dense in consequence.
%
% Notation fixed once in the header rather than page by page. « cof » is his
% own abbreviation for the cofibrations and is kept everywhere, including in
% the subscripts W_cof and C_cof. He writes C_cof for the full subcategory of
% cofibrant objects and cof for the class of arrows, and the two are
% distinguished only by position. « C1 : C3 » is his way of writing a range of
% axioms and is transcribed as he writes it, not expanded to « C1 à C3 ». The
% tilde is doing two different jobs and they must not be confused: over an
% object, x-tilde is a cofibrant replacement or a factorisation object, while
% C-tilde is the new category built out of pi_0 — and over an arrow, a tilde
% written under or beside the shaft is his mark for membership in W, which is
% rendered here by the label « \sim » on the arrow. « f. et s. » is « il faut
% et il suffit ». « comp. conn. » is « composante connexe ». « Ch » is his
% abbreviation and is left as he writes it. His « 1d_x » on page 7 is written
% with a figure one and is kept; « 1_x » elsewhere is the identity.
%
% Readings flagged and not settled: the abbreviation between « cat. de » and
% « co/fibrations » in the corner note of page 1; the subscript under the W of
% the very first line, a small mark that will not resolve and that no later
% occurrence of W repeats; the word after « De ces deux foncteurs » on page 3,
% read « réciproques » on the sense and not on the ink; the noun after « avec
% les » in the N.B. of page 5; the last factor of the composition law at the
% foot of page 5, read « k_y »; the three words after « ce qui » on page 11;
% and almost the whole of the marginal note running up the left edge of
% page 13.
%
% Two variances of substance are flagged where they occur and not repaired.
% The composition law at the foot of page 5 writes its right-hand side
% (g'g)·k_y, in which f does not appear although the left-hand side composes
% (f,i) with (g,j). And page 11's rectified definition of Ch_0 asks for
% f in cof and i in W, where page 5 had asked for i in cof inter W; the two
% definitions stand side by side in the folder under the same name and the
% edition does not reconcile them.
%
% Pass: Opus 5 (claude-opus-5), 2026-09-13 — first pass, unchecked against the pages by a human.
\documentclass[11pt,a4paper]{article}
% Le préambule commun à toutes les transcriptions du fonds Grothendieck.
%
% Il tient en une idée : **l'incertitude se compose, elle ne se résout pas.**
% Une transcription qui lisse un mot douteux a détruit la seule chose pour
% laquelle on la faisait — un lecteur doit pouvoir savoir, à chaque endroit,
% ce qui était sur le feuillet et ce qui a été ajouté. Les six macros qui
% suivent sont donc l'appareil critique, et non de la décoration.
%
% Les mêmes commandes sont reconnues par `scripts/render.mjs`, qui produit la
% vue de lecture du site. Une macro ajoutée ici doit l'être aussi là-bas,
% faute de quoi le rendu échouera bruyamment — ce qui est voulu : un rendu
% silencieusement faux est pire qu'un rendu absent.

\ProvidesPackage{grothendieck}[2026/08/08 Transcriptions du fonds Grothendieck]

\usepackage{amsmath,amssymb,amsthm}
\usepackage{mathrsfs}
\usepackage{stmaryrd}
% Les diagrammes commutatifs sont partout dans ce fonds : c'est la langue
% même de la géométrie algébrique de Grothendieck, et une transcription qui
% les rendrait en prose ne transcrirait rien.
\usepackage{tikz-cd}
% Français par défaut — c'est la langue des feuillets — mais l'anglais est
% chargé aussi : la traduction et le résumé sont des documents anglais, et la
% typographie française (espaces avant les deux-points) y serait une faute.
% Ces documents basculent par \selectlanguage{english} en tête de corps.
\usepackage[english,french]{babel}
\usepackage{fontspec}
\usepackage{geometry}
\geometry{a4paper,margin=25mm,marginparwidth=28mm}
\usepackage{marginnote}
\usepackage{xcolor}
% ulem, not soul. `soul`'s \st and \ul are built on a character-by-character
% scan that XeLaTeX's font machinery breaks: the document fails with an
% undefined control sequence rather than degrading. ulem does the same two jobs
% and is Unicode-engine safe. `normalem` stops it hijacking \emph, which the
% transcriptions use for Grothendieck's own emphasis.
\usepackage[normalem]{ulem}
\usepackage{hyperref}
\hypersetup{colorlinks=true,linkcolor=black,urlcolor=black,pdfborder={0 0 0}}

% --- Métadonnées du lot -----------------------------------------------------
% Reprises telles quelles par le script de rendu, qui les affiche en tête de
% la vue de lecture. Elles fixent ce que le fichier prétend couvrir : sans
% elles, une transcription flotte sans point d'attache dans un fonds de seize
% mille feuillets.

\newcommand{\folder}[1]{\gdef\grothFolder{#1}}
\newcommand{\batch}[1]{\gdef\grothBatch{#1}}
\newcommand{\leaves}[2]{\gdef\grothFirst{#1}\gdef\grothLast{#2}}
\newcommand{\foldertitle}[1]{\gdef\grothFoldertitle{#1}}
\gdef\grothFolder{?}\gdef\grothBatch{?}\gdef\grothFirst{?}\gdef\grothLast{?}\gdef\grothFoldertitle{}

% --- Le résumé --------------------------------------------------------------
%
% Il ouvre la lecture modernisée, sous le titre « Résumé », et il est composé
% en petit. Ce n'est pas un ornement : c'est le seul passage du document qui
% ne soit pas des mathématiques, et un lecteur doit pouvoir le reconnaître
% avant de l'avoir lu — puis le sauter s'il connaît déjà le sujet, ou s'y
% arrêter s'il ne le connaît pas. Le filet qui le suit marque la reprise du
% corps.

\newenvironment{resume}
  {\small\setlength{\parskip}{0.45em}}
  {\par\medskip\hrule\medskip}

% --- Mots-clés ---------------------------------------------------------------
%
% En anglais, en clôture du résumé de la lecture modernisée : le vocabulaire
% moderne sous lequel chercher ce que les feuillets construisent. Le manifeste
% du site (`npm run manifest`) les extrait pour étiqueter la cote entière —
% cette ligne est la source unique des tags, il n'y a pas de fichier à côté.
\newcommand{\keywords}[1]{\par\smallskip\noindent
  {\footnotesize\textsc{Keywords} --- \textit{#1}}\par}

% --- L'appareil critique ----------------------------------------------------

% Le numéro de feuillet, en marge. C'est l'ancre qui permet de régler un
% désaccord sur une formule en regardant *un* feuillet plutôt qu'une chemise
% de six cents. Le site s'en sert aussi pour tourner les pages du fac-similé
% quand on fait défiler la transcription.
\newcommand{\leaf}[1]{%
  \marginnote{\footnotesize\color{gray!70}\textsf{#1}}%
  \label{leaf:#1}\ignorespaces}

% Illisible : on ne devine pas. Un mot inventé qui se lit comme les autres est
% le pire résultat possible.
\newcommand{\ill}{\textcolor{red!60!black}{[\,\ldots\,]}}

% Lecture douteuse : le mot est donné, et signalé comme incertain.
\newcommand{\uncertain}[1]{\uline{#1}}

% Ajout de l'éditeur — une lettre manquante, un mot sous-entendu.
\newcommand{\add}[1]{\textcolor{blue!50!black}{[#1]}}

% Biffé par Grothendieck lui-même. Ce qu'il a rayé fait partie du document :
% une rature dit souvent où la pensée a changé de direction.
\newcommand{\struck}[1]{\sout{#1}}

% Note du transcripteur, sur l'état matériel du feuillet.
\newcommand{\note}[1]{\par\smallskip{\small\color{gray!60!black}\textit{[#1]}}\par\smallskip}

% Note marginale de Grothendieck, distincte du corps du texte.
\newcommand{\marginal}[1]{\par\smallskip\noindent\hspace{1em}%
  {\small\color{gray!60!black}#1}\par\smallskip}

% --- Titre ------------------------------------------------------------------

\AtBeginDocument{%
  \begin{center}
    {\large\bfseries Fonds Alexandre Grothendieck --- cote n\textsuperscript{o}~\grothFolder}\\[2pt]
    {\itshape\grothFoldertitle}\\[4pt]
    {\small feuillets \grothFirst--\grothLast\ (lot \grothBatch)}\\[2pt]
    {\footnotesize\color{gray!60!black}Transcription établie d'après le fac-similé numérisé
     par l'Université de Montpellier.\\
     Les lectures incertaines sont soulignées, les illisibles notées [\,\ldots\,],
     les ajouts entre crochets.}
  \end{center}
  \vspace{1em}}

\endinput


\folder{106}
\batch{1}
\pages{1}{16}
\dating{[à partir de 1982]}
\watermark{Édition de démonstration}
\foldertitle{[Champs (stacks) 3] : notes manuscrites (s.d.)}

\begin{document}

\section{Localisation : $W_{\mathrm{cof}}^{-1}\mathcal{C}_{\mathrm{cof}}
\longrightarrow W^{-1}\mathcal{C}$}

\note{ce titre est de l'édition. Le dossier ne porte aucun titre de sa main,
et celui de l'inventaire, « [Champs (stacks) 3] », est entre crochets, donc
des archivistes. La suite est d'un seul tenant : huit feuillets de bloc à
spirale, écrits au recto seulement, sous sa pagination 1 à 8, qui sont les
pages 1, 3, 5, 7, 9, 11, 13 et 15 des archivistes ; les pages paires sont le
listing d'ordinateur dont il se sert comme papier de rencontre}

\page{1}\note{sa page 1, ouverte sans numéro ; les suivantes portent 2) à 8)}

\marginal{Ho $\mathcal{C}$ \quad cat. de \ill{} co/fibrations ?}

\note{la note est écrite dans le coin supérieur gauche, en travers, et coupée
par un trait réglé en diagonale ; l'abréviation du milieu ne se résout pas}

Soit $(\mathcal{C}, W, \mathrm{cof})$ une catégorie satisfaisant à
C1 : C3 de Baues.

\note{sous le $W$ de cette première ligne il y a une petite marque en indice
qu'aucune des occurrences suivantes ne reprend, et qui ne se lit pas}

\struck{On se place \ill{} qu'il y a \ill{} objet initial.}

(Alors il est connu que

\[
W_{\mathrm{cof}}^{-1}\mathcal{C}_{\mathrm{cof}} \longrightarrow W^{-1}\mathcal{C}
\]

est une équivalence de catégories.)

Pour le voir, on se définit un foncteur en sens inverse (qui sera
quasi-inverse)

\[
\Phi : W^{-1}\mathcal{C} \longrightarrow
W_{\mathrm{cof}}^{-1}\mathcal{C}_{\mathrm{cof}},
\qquad \text{i.e.\ } \mathcal{C} \longrightarrow
W_{\mathrm{cof}}^{-1}\mathcal{C}_{\mathrm{cof}}
\]

qui envoie les $u \in W$ en isos.

Pour tt $x \in \mathcal{C}$, utilisons C3 pour $\emptyset \to x$ : on trouve

\begin{tikzcd}[column sep=small, row sep=small, nodes={font=\scriptsize}]
\widetilde{x} \arrow[r, "\alpha_x"] & x
\end{tikzcd}

avec $\widetilde{x} \in \mathcal{C}_{\mathrm{cof}}$, $\alpha_x \in W$.

Soit alors $x \overset{f}{\longrightarrow} y$ dans $\mathcal{C}$,

\begin{tikzcd}[column sep=small, row sep=small, nodes={font=\scriptsize}]
\widetilde{x} \arrow[d, "\alpha_x"'] \arrow[dr, "f\alpha_x"] & \widetilde{y}
\arrow[d, "\alpha_y"] \\
x \arrow[r, "f"'] & y
\end{tikzcd}

appliquons le « weak lifting lemma », on trouve un relèvement $\widetilde{f}$
de $f\alpha_x$ non \struck{vers $\widetilde{y}$}, mais vers un
$\widetilde{y}\,'$ :

\begin{tikzcd}[column sep=small, row sep=small, nodes={font=\scriptsize}]
\widetilde{x} \arrow[r, "\widetilde{f}"] \arrow[d, no head] &
\widetilde{y}\,' \arrow[d, "\alpha'_y"] \\
\widetilde{x} \arrow[r, "f\alpha_x"'] & y
\end{tikzcd}

\note{à gauche du même feuillet, le triangle qui porte $i$ :
$\widetilde{y} \overset{i}{\to} \widetilde{y}\,'$ dans $W$, avec $\alpha_y$
descendant de $\widetilde{y}$ sur $y$ et $\alpha'_y$ de $\widetilde{y}\,'$ sur
$y$, de sorte que $\alpha_y = \alpha'_y\, i$}

Considérons alors

\[
[i]^{-1}\widetilde{f} : \widetilde{x} \longrightarrow \widetilde{y}
\qquad \text{dans } \mathrm{Hom}_{W_{\mathrm{cof}}^{-1}
\mathcal{C}_{\mathrm{cof}}}(\widetilde{x}, \widetilde{y}).
\]

Je dis qu'il ne dépend pas des choix de $(i, \alpha'_y, \varphi)$. Pour
\ill{} : introduire $\widetilde{y}\,'_3 = \widetilde{y}\,'_1
\vee_{\widetilde{y}} \widetilde{y}\,'_2$, \ill{} deux relèvements \ill{} et on
est ramené (écrivant $y'$ au lieu de $\widetilde{y}\,'$, mais $x$ cof.) à
prouver que si $x', y' \in \mathcal{C}_{\mathrm{cof}}$ et
$\alpha\varphi = \alpha\psi$ \struck{$\alpha g$}, alors $[\varphi] = [\psi]$
dans $W_{\mathrm{cof}}^{-1}\mathcal{C}_{\mathrm{cof}}$ :

\begin{tikzcd}[column sep=small, row sep=small, nodes={font=\scriptsize}]
x' \arrow[r, bend left=20, "\varphi"] \arrow[r, bend right=20, "\psi"'] & y'
\arrow[r, "\underset{\sim}{\alpha}"] & y
\end{tikzcd}

\note{plus bas, la même chose écrite par la somme : $\nabla_{x'}$ descend de
$x' \sqcup x'$ sur $x'$, et $(\varphi, \psi)$ va de $x' \sqcup x'$ dans $y'$.
Un bloc annulé de trois lignes, à droite, reprenait $x \to y$ avec les mêmes
hypothèses de cofibration}

Ceci ayant lieu, on trouve un foncteur $\mathcal{C} \longrightarrow
W_{\mathrm{cof}}^{-1}\mathcal{C}_{\mathrm{cof}}$, la compatibilité aux
identités et aux compositions étant immédiate.

\page{2}

\note{feuillet de listing tourné d'un quart de tour et couvert, sur toute sa
largeur, d'un champ de petits diagrammes d'essai : ce sont, à quelques
variantes près, toujours le même — la comparaison de deux relèvements de
$f\alpha_x$ dans deux objets $\widetilde{y}\,'$ et $\widetilde{y}\,''$, et la
somme amalgamée des deux sous $\widetilde{y}$. Presque aucun n'est étiqueté et
il n'y a pas une ligne de prose sur le feuillet. Les redessiner un à un
reviendrait à leur prêter des étiquettes que la page ne porte pas ; ils ne
sont donc pas transcrits, et c'est la page 1 qui donne la figure sous sa forme
écrite}

\page{3}\note{sa page 2}

\note{le tiers supérieur du feuillet porte deux bandes de diagrammes d'essai,
plusieurs fois repris et en partie biffés. Le seul qui se lise est un prisme à
trois étages : $\widetilde{x} \to \widetilde{y} \to \widetilde{z}$ en haut,
une rangée intermédiaire $\widetilde{x} \to \widetilde{y}\,' \to
\widetilde{z}\,''$, et $x \to y \to z$ en bas, les verticales étant les
$\alpha$ et portant le signe $\sim$}

Montrons que ce foncteur transforme $f \in W$ en iso. En effet, si $f \in W$,
donc $f\alpha_x \in W$, on voit aussitôt que \ill{} $\widetilde{f} \in W$,
donc $[\widetilde{f}]$ est un iso.

Considérons les deux composés

\begin{tikzcd}[column sep=small, row sep=small, nodes={font=\scriptsize}]
W_{\mathrm{cof}}^{-1}\mathcal{C}_{\mathrm{cof}} \arrow[r, bend left=20, "\Psi"]
& W^{-1}\mathcal{C} \arrow[l, bend left=20, "\Phi"]
\end{tikzcd}

\struck{prouvons que ces deux} De ces deux foncteurs
\uncertain{réciproques} \ill{}.

Pour $\Phi\Psi$, c'est $=$ l'identité ; \ill{} choisir
$\widetilde{x} \overset{\alpha_x}{\to} x$ \struck{\ill{}} \ill{} quand
$x \in \mathcal{C}_{\mathrm{cof}}$. Pour $\Psi\Phi$, il \struck{suffit}
\add{faut} voir que le foncteur \struck{$\Psi\Phi|\mathcal{C}$}

\[
R \overset{\text{déf}}{=} \Psi\Phi|\mathcal{C},
\qquad x \longmapsto \widetilde{x},
\qquad \mathcal{C} \longrightarrow W^{-1}\mathcal{C}
\]

est \struck{bien} isomorphe au foncteur canonique. On a

\[
R(x) = \widetilde{x} \overset{[\alpha_x]}{\longrightarrow} x
\]

dans $\mathcal{C}$ \struck{ce} donc dans $W^{-1}\mathcal{C}$, qui est un iso ;
il faut \ill{} vérifier que c'est fonctoriel, i.e.\ qu'on a

\[
[\alpha_y]\,\bigl([i_f]^{-1}[\widetilde{f}]\bigr) = [f]\,[\alpha_x].
\]

\begin{tikzcd}[column sep=small, row sep=small, nodes={font=\scriptsize}]
\widetilde{x} \arrow[r, "\widetilde{f}"] \arrow[d, "\underset{\sim}{\alpha_x}"'] &
\widetilde{y} \arrow[r, "i_f"] \arrow[d, "\underset{\sim}{\alpha_y}"] &
\widetilde{y}\,' \arrow[dl, "\beta_f"] \\
x \arrow[r, "f"'] & y &
\end{tikzcd}

Mais $\alpha_y = \beta_f\, i_f$, d'où le premier \ill{}, puisque
\struck{\ill{}} on a

\[
[\beta_f \widetilde{f}] = [f\alpha_x] = [f][\alpha_x]
\qquad \text{cqfd.}
\]

\page{5}\note{sa page 3}

En résumé, on a trouvé

\begin{quote}
\textbf{Prop.} Pour que
$\Psi : W_{\mathrm{cof}}^{-1}\mathcal{C}_{\mathrm{cof}} \longrightarrow
W^{-1}\mathcal{C}$ soit une équivalence de cat., il f.\ et s.\ \add{?} que la
condition suivante soit vérifiée :

($\star$) Pour tout diagramme

\begin{tikzcd}[column sep=small, row sep=small, nodes={font=\scriptsize}]
x \arrow[r, bend left=20, "\varphi"] \arrow[r, bend right=20, "\psi"'] & y
\arrow[r, "\underset{\sim}{f}"] & y'
\end{tikzcd}

\add{avec} $f\varphi = f\psi$, $f \in W$, $x, y \in
\mathcal{C}_{\mathrm{cof}}$, on a $[\varphi] = [\psi]$ dans
$W_{\mathrm{cof}}^{-1}\mathcal{C}_{\mathrm{cof}}$ (donc $[\varphi] = [\psi]$
dans $W^{-1}\mathcal{C}$).
\end{quote}

(condition impliquée par celle que $\Psi$ soit \emph{fidèle})

Nous allons travailler maintenant avec $\mathcal{C}_{\mathrm{cof}}$ au lieu de
$\mathcal{C}$ — \struck{sans} quitte à changer les notations, i.e.\ on suppose
que tt objet de $\mathcal{C}$ est cofibrant.

[ N.B.\ $\mathcal{C}_{\mathrm{cof}}$, avec les \ill{} induites par $W$ et cof,
est encore une catégorie \struck{satisfaisant à} C1 : C3, littéralement,
\ill{} de $\mathcal{C}_{\mathrm{cof}}$ dans $\mathcal{C}$ commutent aux sommes
amalgamées \struck{qui} relatives : des diagrammes

\begin{tikzcd}[column sep=small, row sep=small, nodes={font=\scriptsize}]
a \arrow[r, "f"] \arrow[d, "g"'] & b \\
c &
\end{tikzcd}

telles que $f$ ou $g$ soit dans cof (les seules sommes dont on postule
l'existence !). ]

Soit, pour $x, y$ dans $\mathcal{C}$,

\[
\underline{Ch}_0(x,y) = \text{catégorie des diagrammes}
\]

\begin{tikzcd}[column sep=small, row sep=small, nodes={font=\scriptsize}]
x \arrow[r, "f"] & \widetilde{y} & y \arrow[l, "\underset{\sim}{i}"']
\end{tikzcd}

où $i \in W \cap \mathrm{cof} = \mathrm{cof} \cap W$, \struck{catégorie}
(petite, connexe \ill{}). On définit des accouplements

\[
\underline{Ch}_0(x,y) \times \underline{Ch}_0(y,z) \longrightarrow
\underline{Ch}_0(x,z),
\qquad (F, G) \longmapsto G \cdot F
\]

\begin{tikzcd}[column sep=small, row sep=small, nodes={font=\scriptsize}]
& \widetilde{y} \arrow[r, "g"] & \widetilde{z} \arrow[d, "\underset{\sim}{k}"] & z
\arrow[l, "\underset{\sim}{j}"'] \\
& & \widetilde{z}\,' &
\end{tikzcd}

\note{la flèche $\widetilde{z} \to \widetilde{z}\,'$ est marquée cocart sur le
feuillet, et $g'$ désigne le composé $\widetilde{y} \to \widetilde{z}\,'$}

\[
(g,j) \circ (f,i) = (g'g) \cdot \uncertain{k_y}
\]

\note{le membre de droite, tel qu'il est écrit, ne fait pas intervenir $f$,
alors que le membre de gauche compose $(f,i)$ avec $(g,j)$ ; le dernier
facteur est lu $k_y$ et n'est pas assuré. La variance est signalée et non
corrigée}

\section{La catégorie $\widetilde{\mathcal{C}}$}

\note{ce titre est de l'édition ; la page 7 nomme la catégorie mais ne porte
pas de titre}

\page{7}\note{sa page 4}

Cette composition est associative à iso can.\ près, et elle a des identités
bilatères (à un iso can.\ !) \ill{}

\begin{tikzcd}[column sep=small, row sep=small, nodes={font=\scriptsize}]
x \arrow[r, "1"] & \widetilde{x} & x \arrow[l, "\underset{\sim}{1}"']
\end{tikzcd}

noté $\underline{1}_x$.

\struck{On} Définissons \struck{par des \ill{}} des foncteurs

\[
\underline{\mathrm{Hom}}_{\mathcal{C}}(x,y)_{\text{discret}}
\longrightarrow \underline{Ch}_0(x,y),
\qquad f \longmapsto (f, 1d_x)
\]

\note{« $1d_x$ » est écrit avec un chiffre un et est conservé tel quel}

\struck{compatibles avec} compatibles avec les compositions. Ceci posé, soit

\[
\pi_0(x,y) \overset{\text{déf}}{=} \pi_0\bigl(\underline{Ch}_0(x,y)\bigr),
\]

alors le composé dans $\underline{Ch}_0$ définit une composition des
$\pi_0(x,y)$, qui est associative \ill{} et admet les $\underline{1}_x$ comme
unités bilatères. D'où une catégorie, que je note $\widetilde{\mathcal{C}}$,
et un foncteur

\[
\mathcal{C} \longrightarrow \widetilde{\mathcal{C}}.
\]

Pour l'instant, on n'a utilisé \struck{rien} que $(\mathcal{C}, W)$,
\struck{sauf} \ill{} $\mathcal{C}$ donnant des cof.\ $\subset
\mathrm{Fl}\,\mathcal{C}$, \struck{on ne convient de prendre les} \ill{}
\add{(Conv.\ $W$, \struck{cofibs})} — mais \ill{} peut prendre la
construction. Mais $\mathrm{cof} \cap W = \struck{\ill{}}$ \ill{}. On \ill{}
foncteur transforme $f \in W$ en iso : on cherche $g, i$ tels que

\[
(gf, i) \in \text{comp.\ conn.\ de } \underline{1}_x
\text{ dans } \pi_0\bigl(\underline{Ch}_0(x,x)\bigr)
\]

et $(\overline{f}g, j) \in$ comp.\ conn.\ de $\underline{1}_y$.

\note{deux figures accompagnent ce dernier passage : $x \to y$ avec $g$
remontant, et un carré $x \to \widetilde{y}$ au-dessus de $\bar{x} \to \bar{y}$
dont la flèche du bas est marquée cocart, les verticales portant $i$ et le
signe $\sim$}

\page{9}\note{sa page 5}

On utilise le weak lifting lemme :

\begin{tikzcd}[column sep=small, row sep=small, nodes={font=\scriptsize}]
\emptyset \arrow[r] \arrow[d] & x \arrow[r, "\underset{\sim}{i}"] \arrow[d, "\underset{\sim}{f}"'] & \bar{x} \arrow[dl, "\alpha", dashed] \\
y \arrow[r, "1"'] & y \arrow[ur, "g"', dashed] &
\end{tikzcd}

donc il existe $i, \alpha$ \ill{} $\alpha i = f$, \struck{\ill{}}
$i \in W \cap \mathrm{cof}$, et $\alpha$ admet une section $g$,

\[
\alpha g = 1_y.
\]

\marginal{N.B.\ On utilise $y \in \mathcal{C}_{\mathrm{cof}}$
\struck{\ill{}} ici}

Je dis que $(g,i)$ est l'inverse cherché. \struck{Considérons alors \ill{}}

Pour $(gf, i) \sim \underline{1}_x$, on a envie d'utiliser le diagramme
commutatif ci-contre, le seul ennui c'est que, comme $\alpha \in W$ — mais pas
\ill{} $\alpha \in \mathrm{cof}$, donc $\alpha i \ (= f)$ ne vit pas dans
\struck{cof} (alors qu'il est dans $W$).

\begin{tikzcd}[column sep=small, row sep=small, nodes={font=\scriptsize}]
x \arrow[r, "\underset{\sim}{i}"] \arrow[d, "1_x"'] \arrow[dr, "gf"] & \bar{x}
\arrow[d, "\underset{\sim}{\alpha}"] \\
x \arrow[r, "f"'] & y
\end{tikzcd}

Ça marcherait si on pouvait trouver $y \overset{\beta}{\longrightarrow} y'$
\add{dans $W$} tel que le composé $\beta f$ soit non seulement dans $W$, mais
en plus, dans cof. Visiblement, ce n'est pas \struck{ça} le cas (prendre le
cas où les cof.\ sont les monos, et où $f$ n'est pas \uncertain{un mono} !).
Ça se présente mal !

On a donc plutôt \struck{définit} définir d'abord le foncteur canonique

\[
q : \widetilde{\mathcal{C}} \longrightarrow W^{-1}\mathcal{C},
\qquad x \longmapsto x,
\qquad (f,i) \longmapsto [i]^{-1}[f]
\]

il est immédiat que c'est compatible avec les identités et le composition.
Donc on a \struck{un} des foncteurs \struck{surjectif} \add{bijectif} sur les
objets. Je dis que c'est \ill{} surjectif sur les flèches ; il suffit de voir,
pour \ill{} $f : x \to y$ \ill{} dans $W$, que $[f]^{-1}$ provient de
$\widetilde{\mathcal{C}}$, défini par $(g,i)$. À priori \struck{\ill{}}
\struck{$(f,\star)$} $y \to x$ dans $\widetilde{\mathcal{C}}$, défini par
$[i]^{-1}[g]$ ; pour prouver que $[i]^{-1}[g] = [f]^{-1}$, il suffit de
prouver que

\[
[f]\,[i]^{-1}[g] = [1_y],
\]

or $\alpha i = f$ implique $[\alpha][i] = [f]$, donc $[f][i]^{-1} = [\alpha]$,
\ill{} les produits, car $[\alpha][g] = [\alpha g] = [1_y]$, cqfd.

\page{11}\note{sa page 6}

La difficulté, c'est qu'il n'est pas clair que le foncteur

\[
\widetilde{\mathcal{C}} \longrightarrow W^{-1}\mathcal{C}
\]

précédent soit aussi \emph{fidèle} \add{$\alpha = $ un iso} (ce qui
impliquerait que c'est une équivalence de catégories).
\struck{Dans ce cas \ill{} le foncteur \ill{} viendrait \ill{}.}

\struck{\ill{}} que le foncteur canonique $\mathcal{C} \longrightarrow
\widetilde{\mathcal{C}}$ (dont le composé avec le foncteur can.\
$\mathcal{C} \to W^{-1}\mathcal{C}$) \ill{} bien \ill{} $W$ en iso, on \ill{}
que \ill{} transforme les flèches de $W$ \ill{} constructions qui \ill{}
$[f, i]$ de $\widetilde{\mathcal{C}}$, est bien utile, \ill{} d'ailleurs
\struck{\ill{}}, \add{semble-t-il}, de poser

\[
W_0 = W \cap \mathrm{cof}.
\]

Cela est quinait d'ailleurs \ill{} le foncteur

$W_0^{-1}\mathcal{C} \longrightarrow$ \ill{}

devrait une équivalence de catégories, \ill{} le foncteur
$\widetilde{\mathcal{C}} \to W^{-1}\mathcal{C}$ se factorise visiblement par

\begin{tikzcd}[column sep=small, row sep=small, nodes={font=\scriptsize}]
\widetilde{\mathcal{C}} \arrow[r] & W_0^{-1}\mathcal{C} \arrow[r] &
W^{-1}\mathcal{C}
\end{tikzcd}

Mais je doute qu'il ne suffise d'inverser les équiv.\ faibles
\emph{cofibrantes}, pour que les équiv.\ faibles deviennent inversibles !

Je rectifie donc la définition, en prenant plutôt la catégorie (notre
précédent \struck{\ill{}} $\widetilde{\mathcal{C}}$) \ill{} cette catégorie
\struck{définie via les flèches} $\pi_0$ des $\underline{Ch}_0(x,y)$,
\struck{définies \ill{}} (pour $x, y$ donnés) comme \ill{} celles des
diagrammes

\begin{tikzcd}[column sep=small, row sep=small, nodes={font=\scriptsize}]
x \arrow[r, "f"] & \widetilde{y} & y \arrow[l, "\underset{\sim}{i}"']
\end{tikzcd}

avec $\begin{cases} f \in \mathrm{cof} \\ i \in W \end{cases}$

\note{c'est le point de bascule du dossier : la page 5 demandait
$i \in \mathrm{cof} \cap W$ sans rien demander à $f$, et la définition
rectifiée demande $f \in \mathrm{cof}$ et $i \in W$. Les deux définitions
portent le même nom $\underline{Ch}_0$ et se suivent dans le dossier ;
l'édition ne les concilie pas}

\page{12}

\note{feuillet de listing, daté 24 AUG 82 par la machine, dont la marge
inférieure porte deux petites figures d'essai de sa main : des objets tous
notés $x$, $x'$, $x''$, reliés par des flèches montantes dont l'une est
biffée. Aucune n'est étiquetée et elles ne portent pas de mathématique qu'on
puisse restituer ; elles ne sont donc pas redessinées. La date imprimée est
relevée ici parce que c'est elle, et celle de la page 14, qui datent le
dossier}

\page{13}\note{sa page 7}

\note{le feuillet le moins lisible du dossier : un palimpseste rapide, deux
blocs annulés par de longs traits horizontaux, et une note marginale écrite en
biais le long du bord gauche que la résolution du film ne permet pas de
restituer}

\ill{} que le composé \ill{} comme précédemment. \ill{}

\struck{$\widetilde{\mathcal{C}} \longrightarrow W^{-1}\mathcal{C}$}
\struck{\ill{} foncteurs transforme bien} \ill{} flèches inversible.
\emph{On se} propose \ill{}

\[
F = [f], \qquad G = [1_y, f]
\]

\[
\mathcal{C} \longrightarrow \widetilde{\mathcal{C}}
\]

qui ne \ill{} pas plus \ill{} devrait une factorisation \ill{}
$\widetilde{\mathcal{C}}$ \ill{} plus délicat. Soit $f : x \to y$ \ill{}

\begin{tikzcd}[column sep=small, row sep=small, nodes={font=\scriptsize}]
x \arrow[r, "f_0"] & \bar{y} \arrow[r, "\underset{\sim}{j}"] & \widetilde{y}\,'
\end{tikzcd}

\ill{} (et $y$ étant cofibrant) \ill{} les précédentes, où $\beta$ a une
section $i$.

$[f_0, i]$ : il faut montrer \ill{} deux faits. Pour on écrivait $[f, i]$
\ill{} qu'il existe \struck{\ill{}} une rétraction $\alpha$ de
$\widetilde{y}$ sur $y$, telle que

\[
\alpha f' = f
\]

\ill{}, il faut voir qu'ils sont \ill{} dans \ill{}
$\underline{Ch}_0(x,y)$. Mais pour deux factorisations $\alpha f_0$ et
$\alpha' f'_0$, avec $f_0, f'_0 \in \mathrm{cof}$, $\alpha, \alpha' \in W$,
\ill{} une troisième \struck{\ill{}} $\alpha'', f''_0$ qui \ill{} soit dans le
diagramme \ill{}

\page{14}

\note{feuillet de listing, daté 25 AUG 82 par la machine, dont la marge droite
porte, de sa main, le bloc de sommes amalgamées que la page 15 met en œuvre.
Les trois carrés se lisent, le carré du milieu étant marqué cocart, $i_1$ et
$i_2$ étant les deux insertions}

\begin{tikzcd}[column sep=small, row sep=small, nodes={font=\scriptsize}]
x \arrow[r] \arrow[d, "i_1"'] & y' \arrow[d] \\
x \vee x \arrow[r] \arrow[d] & y' \vee x \arrow[d] \\
x \vee y'' \arrow[r] & y' \vee y''
\end{tikzcd}

\note{au-dessus du bloc, deux figures reprises : $x \to y'$ avec $y' \to y$ et
$y'' \to y'$ ; puis $x \amalg x \to y' \amalg y''$, suivi d'un mot biffé. À
droite du carré supérieur, $x \to \bullet \to y$ avec $x \to y'$ en dessous,
$i_2$ portant l'insertion de gauche}

\page{15}\note{sa page 8}

\begin{tikzcd}[column sep=small, row sep=small, nodes={font=\scriptsize}]
x \vee x \arrow[r, "(f_0) \vee (f'_0)"] \arrow[d, "\nabla_x"'] & \bar{y}
\amalg \bar{y}\,' \arrow[d, dashed] \arrow[dr, "{(\alpha, \alpha')}"] & \\
x \arrow[r, "g"', dashed] & \widetilde{y} \arrow[r, "\varphi_0"'] & \bar{y}\,''
\arrow[d, "\alpha''"] \\
& & y
\end{tikzcd}

\note{le carré de gauche est marqué cocart ; sous la figure, une accolade
ramène le tout à $f$}

on factorise $\varphi$ \ill{} en $\alpha''\varphi_0$ \ill{}
($\varphi_0 \in \mathrm{cof}$, $\alpha'' \in W$) et on prend

\[
f''_0 = \varphi_0\, g.
\]

\ill{} des sections, il suffirait que $\alpha$ et $\alpha'$ les aient, \ill{}
$\alpha''$ (qui leur \ill{}) \ill{}. On a \ill{} une flèche

\[
u : (f_0, \alpha) \longrightarrow (f''_0, \alpha'')
\]

dans le \ill{} objets \ill{}.

\begin{tikzcd}[column sep=small, row sep=small, nodes={font=\scriptsize}]
x \arrow[rr, "f"] \arrow[dr, "f_0"'] \arrow[ddr, "f'_0"'] & & y \\
& \bar{y} \arrow[ur, "\alpha'"'] \arrow[d, "u"'] & \\
& \bar{y}\,'' \arrow[uur, "\alpha''"'] &
\end{tikzcd}

Donc \emph{si} on choisit une \struck{le choix de} section $i'$ \ill{} comme
l'image par \ill{} de ce dernier $i'$ de $\alpha'_0$ \ill{}, on \ill{}. Donc
il reste à prouver \ill{}.

Alors \ill{} choisir de $f$ \ill{} $\alpha f_0$, $[f_0, i]$ ne dépend pas des
choix, comme

\[
[f_0, i] = [\mathrm{id}_{\bar{y}}, i] \circ (f, \mathrm{id}_{\bar{y}}),
\]

\ill{} à ceci :

\begin{quote}
($\star\star$) Soit $x \overset{f}{\longrightarrow} y$ dans $W$, tel que $f$
ait \add{deux $i, i' \in W$} \struck{deux \ill{}, sections} $i, i'$ ; alors
\struck{le \ill{} \ill{} de} la classe dans
$\pi_0\bigl(\underline{Ch}_0(x,y)\bigr)$ de $(1_x, i)$ ne dépend pas des choix
de la section $i$.
\end{quote}

\note{à gauche, la figure qui l'accompagne : $x \overset{f}{\to} \bar{y}$ avec
$\alpha$ et les deux arcs $i$ et $i'$ remontant de $\bar{y}$ sur $x$ ; puis
$x \overset{1_x}{\to} x$ avec $y$ au-dessus et $i$ descendant}

Cette condition est (nec.\ et) suffis.\ pour que $\widetilde{\mathcal{C}} \to
W^{-1}\mathcal{C}$ soit \emph{fidèle} (donc un iso), comme on \ill{}
précis\ill{} \ill{} une construction \ill{} tel \ill{} et bien le foncteur
canonique

\[
\mathcal{C} \longrightarrow \widetilde{\mathcal{C}}.
\]

\end{document}
